% Proposed insertion after MN.1's normal-divisor table. This supplies a
% reproducible local certificate and does not replace the full CFL/IL gates.
\paragraph{An explicit reconstruction of the eleven native charts.}
The moving-star input can be recovered directly from the raw zero cones.
In the original ordered lattice $(X,Y,P,p)$, retain $A,B_1,B_2,B_3$
of NS.3 and put
\[
 U_1=(-5,0,1,0),\qquad U_2=(-10,-3,1,0),\qquad
 U_3=(-13,-6,1,0).
\]
The six raw cones meeting the two physical faces are
\[
 \langle A,B_1,U_1,U_2\rangle,\quad
 \langle A,B_1,(-10-k,-3-k,1,0),(-11-k,-4-k,1,0)\rangle
 \quad(0\le k\le2),
\]
\[
 \langle A,B_1,U_3,B_2\rangle,\qquad
 \langle A,B_3,U_3,B_2\rangle.
\]
Insert the two specified I2Z.14 rays $(-9,-2,1,1)$ and
$(-7,-1,1,1)$ into every containing cone. Then, for $k=1,2,3$,
successively insert $(-10-k,-3-k,1,0)+jA$, $j=1,\ldots,k$.
Each latter insertion is the sum of the two then-existing rays
in its source-cut edge, hence has the two ordinary blowup maps
$(u,v)\mapsto(u,uv)$ and $(u',v')\mapsto(u'v',v')$.
Retaining every child gives a refinement with the same support.

The seven charts containing $\langle A,B_1\rangle$ are obtained
by adjoining consecutive pairs of the following ordered list:
\[
\begin{split}
 &(-5,0,1,0),\ (-7,-1,1,1),\ (-9,-2,1,1),\
 &(-10,-3,1,0),\ (-11,-4,1,1),\ (-12,-5,1,2),\
 &(-13,-6,1,3),\ B_2.
\end{split}
\]
The four charts containing $\langle B_2,B_3\rangle$ are obtained
by adjoining consecutive pairs of
\[
 A,\quad U_3+3A,\quad U_3+2A,\quad U_3+A,\quad U_3.
\]
All eleven ray matrices are unimodular. Each chart lies entirely
on one side of the actual source functional $y+3r+z=0$.
Each has at most one upper ray $z=0$, so the upper-face common
lift of I2ZC.2 retains these whole charts. The actual period $T_0$
maps $B_3$ to $A$, $B_2$ to $B_1$, and $U_3$ to $U_1$;
it therefore gives their written closing face, including its
$f_0$ coefficient and reciprocal.

For the first seven charts evaluate the native infinity rows
\[
 m_-^\infty=(1,-1,1,1),\qquad m_+^\infty=(-1,2,0,0).
\]
For the remaining four charts first apply the full $T_0$ lattice
matrix, then evaluate
\[
 m_-^0=(1,-2,1,1),\qquad m_+^0=(-1,1,0,0).
\]
These are exactly NS.23 and the native overlap NS.15. On the
ordered normal rays their values are $(1,0)$ and $(0,1)$;
on the horizontal tangent rays both values are zero. On the ten
compact tangent rays they give, in the order $C_3,\ldots,C_{-6}$,
\[
 (-4,-4,-4,-4,-4,-5,-6,-5,-4,-3),\qquad
 (5,5,5,5,5,5,4,3,2,1).
\]
Thus every entry of the divisor table follows from the original
four-dimensional source and line characters. The resulting
wall-intersection form gives $D_-^2=-5$, $D_-D_+=2$ and
$D_+^2=-4$.

The exact reconstruction program
\path{code/K111_moving_star_reconstruction.py} and its output
\path{K111_moving_star_reconstruction.json} retain the complete
stellar maps, all eleven matrices and inverses, actual raw-parent
invariant-semigroup pullbacks, physical weights, canonical orders,
and both normal derivative rows. They certify this whole moving
star; the compact and lower-reference ledgers have their separately
specified inputs.
